\newpage
\section*{1\quad The problem, and seventy years of attempts}
\addcontentsline{toc}{section}{1 The problem}
\startchapter{1}

\begin{chaptermap}
\emph{The problem:} a customer writes \tok{waste of time}, and we want a machine to write a good reply.
\emph{The steps:} three ways to automate the reply, and why writing it one word at a time wins (\S1.1--1.2);
then seventy years of attempts at the one function that needs, each with the failure that pushed the next
one (\S1.3--1.7); then a real case (\S1.8) and the whole machine at a glance (\S1.9).
\emph{Where it ends:} with four requirements that the transformer was built to meet, and a picture of how
one pass handles a whole batch of sentences.
\end{chaptermap}

\begin{beforeyoustart}
Five quick questions. Answer each, then read the refresher under it.
\begin{enumerate}[label=\arabic*.,leftmargin=1.6em]
\item A bag holds 3 red counters and 1 blue counter. You pick one without looking. What is the chance it is red?\\
\emph{Refresher:} a probability is ``favourable cases divided by all cases'', always between 0 and 1. Here
$3/4 = \nOneRed$.
\item You toss a fair coin twice. What is the chance of two heads?\\
\emph{Refresher:} for two separate events, ``this and then that'' means multiply: $\tfrac12\times\tfrac12=\nOneTwoHeads$.
\item What is $0.9\times0.9\times0.9$?\\
\emph{Refresher:} multiplying a number by itself is written as a power: $0.9^3 = \nOneCube$. A number below 1
gets smaller every time you multiply by it.
\item A friend's calculation says the chance of rain tomorrow is 1.2. What went wrong?\\
\emph{Refresher:} nothing can be more certain than certain (1) or less than impossible (0). Any probability
outside 0 to 1 means a mistake earlier in the working: a good thing to notice, because it tells you where to look.
\item What is 30\% of 200?\\
\emph{Refresher:} ``per cent'' means ``out of 100'', so 30\% of 200 is $0.30\times200=\nOnePct$.
\end{enumerate}
\end{beforeyoustart}

Before any architecture, we need a problem. Nobody wakes up wanting a transformer. We want something done,
and the transformer is what fell out after many people tried many simpler things. So let us start where they
started: with a job.

%---------------------------------------------------------------------
\subsection*{1.1\quad A real job}

A software company sends every customer a message after a support chat: \emph{How was your experience?}
Thousands of answers come back every day. Most say \emph{thanks}. Some say \emph{great}. And some say
\begin{center}\tok{waste of time}\end{center}
What should the company say back? Tell me: if you were the support agent, what would you write?

A good agent writes something like \emph{Sorry, this has not been useful}, and does not argue, does not paste a
refund policy, and does not reply ``Glad we could help!''. Our job: get a machine to write that reply.

How did you, the human, do it? You read the whole message. You noticed that \tok{waste} turns \tok{time} into a
complaint. You wrote the reply one word after another, each word fitting the words before it. Keep those three
things in mind: \emph{read everything}, \emph{notice which words change the meaning of which}, \emph{write one
word at a time}. The whole transformer is a machine for doing these three things.

\Summary{Our running problem is to reply well to \tok{waste of time}, which needs reading everything, noticing
which words change which, and writing one word at a time.}

%---------------------------------------------------------------------
\subsection*{1.2\quad Three ways to automate the reply}

\paragraph{Attempt 1: rules.} \emph{If the message contains ``waste of time'', reply with an apology.}
Will this survive a week? On Monday it works. On Tuesday someone writes \emph{total waste}, \emph{useless},
\emph{not worth my evening}, and the rule misses all three. On Wednesday someone writes \emph{not a waste of
time at all, loved it!} and the rule apologises to a happy customer. Rules break because language has endlessly
many ways to say the same thing, and small words like \emph{not} flip the meaning of everything around them.

\paragraph{Attempt 2: pick a reply from a list.} Write ten good replies and train a model to choose one
(``apology'', ``thanks'', ``escalate'', \ldots). This learns from examples instead of rules, which is better.
But what happens with the eleventh situation, a customer who is angry \emph{and} asks a question? There is no
reply for it in the list. The model can only pick from what we imagined in advance.

\paragraph{Attempt 3: write the reply, one word at a time.} Instead of choosing a finished reply, the machine
builds one. And here is the trick: \emph{writing a reply is guessing the next word, over and over.}

\Picture When a teacher says ``Twinkle, twinkle, little \rule{1.2em}{0.4pt}'', a room of children says
\emph{star}. They are not reciting a rule. They have heard so much language that, given the words so far, one
word is far more likely than all others. Play that game again on \emph{Twinkle, twinkle, little star, how I
\rule{1.2em}{0.4pt}}, and again, and again, and a whole verse comes out, one guess at a time.

\begin{breaks}
The children answer with one word. The model answers with a probability for \emph{every} word it knows, and
something else (Chapter 7) picks one. The children also know what a star is; the model has only ever seen text.
\end{breaks}

\Exactly Let a piece of text be a sequence of tokens $w_1, w_2, \ldots, w_T$ (for now, a token is a word;
Chapter 2 refines this). The \emph{chain rule of probability}, which is ``first this, then that given this'',
says
\[
P(w_1, w_2, \ldots, w_T) \;=\; P(w_1)\,P(w_2\mid w_1)\,P(w_3\mid w_1,w_2)\cdots P(w_T\mid w_1,\ldots,w_{T-1})
\;=\; \prod_{t=1}^{T} P(w_t \mid w_{<t}).
\]
In words: the chance of the whole text is the chance of the first word, times the chance of the second word
given the first, and so on. It is the coin question from \emph{Before you start}, with each step allowed to
depend on the earlier ones. So if a machine answers one question well (\emph{given the words so far, how
likely is each possible next word?}), it can write anything. Our reply
\[
\underbrace{\tok{waste of time}}_{\text{given}} \;\to\; \tok{Sorry} \;\to\; \tok{,} \;\to\; \tok{this}
\;\to\; \tok{has} \;\to\; \tok{not} \;\to\; \tok{been} \;\to\; \tok{useful}
\]
is seven guesses, each made with everything to its left in view. Figure~\ref{fig:loop} shows the loop.

\begin{figure}[h]
\centering
\begin{tikzpicture}[font=\small, box/.style={draw, rounded corners=2pt, inner sep=5pt, align=center}]
\node[box] (ctx) {words so far\\\tok{waste of time Sorry ,}};
\node[box, right=1.1cm of ctx] (model) {the model};
\node[box, right=1.1cm of model] (prob) {a probability for\\every word it knows\\[2pt]
\footnotesize \tok{this}: high,\ \tok{we}: lower,\ \ldots};
\node[box, below=0.7cm of model] (pick) {pick one: \tok{this}};
\draw[-{Stealth}] (ctx) -- (model);
\draw[-{Stealth}] (model) -- (prob);
\draw[-{Stealth}] (prob) |- (pick);
\draw[-{Stealth}] (pick) -| node[pos=0.75, left, font=\footnotesize]{append, repeat} (ctx);
\end{tikzpicture}
\caption{Writing as repeated guessing. Chapter 7 shows how the probabilities are computed. Everything in
Chapters 2--6 lives inside the box called ``the model''.}
\label{fig:loop}
\end{figure}

\Predict Let's assume each of the seven guesses is right with probability 0.9, and that one wrong word spoils
the reply. Before computing: is the chance that the whole reply comes out right closer to 90\%, 70\% or 50\%?

The chain rule says multiply: $0.9^7 = \nOneSevenRight$, about \nOneSevenPct\%. For a 20-word reply it would be
$0.9^{20}=\nOneTwentyRight$.
\Sanity Every factor is below 1, so the product must fall as the reply grows, and it does.

\Why Two consequences fall straight out of this loop, and you should hold on to both.
\begin{enumerate}[label=(\roman*)]
\item \emph{Every word is written in the light of all earlier words, including the model's own.} If one guess
goes wrong, every later guess builds on the mistake, and the worked example shows how fast that adds up. A large
part of what people call \emph{hallucination} comes from here.
\item \emph{The problem labels itself.} Any text is already a pile of solved exercises: hide the next word, and
the text itself is the answer key. A corpus of \cTrainWords\ words of transcribed talks is \cTrainWords\ such
exercises, minus one, and nobody had to label a single one.
\end{enumerate}

\begin{wrongidea}{1.3}
\emph{Hallucination is a bug someone forgot to fix.} It is tempting because software errors usually are bugs.
But look at consequence (i): the same loop that lets the model write a fluent reply to a message it has never
seen also lets one wrong guess carry the rest of the reply with it. Hallucination is a consequence of the process, not a
separate fault sitting in the code. Systems built on these models reduce it with
checks around the model; they do not delete a line of code.
\end{wrongidea}

\begin{wrongidea}{1.4}
\emph{Training data needs a human to label every example.} Every model you may have built before (a churn
model, a fraud model, a release-triage model) needed a labelled $X$ and $Y$, so this feels like a law. For the
next-word game the label is the next word, already sitting in the text. Attempt 2 needed people to label
replies; Attempt 3 needs only text.
\end{wrongidea}

\begin{wrongidea}{1.1}
\emph{Next-word prediction is shallow autocomplete, so it cannot involve understanding.} Phone keyboards also
predict the next word, and they are not clever, so the comparison is natural. But ask what it takes to guess
the next word \emph{well} after \tok{waste of time}: you need to know it is a complaint, who is speaking, and
what a polite company says. The game is simple to state and hard to play; the better a model plays it, the more
it has to capture about the world the text describes.
\end{wrongidea}

\begin{wrongidea}{1.2}
\emph{When a model says ``I am 90\% sure'', that is a calibrated probability.} It looks like one, and the
model does compute probabilities inside. But those are probabilities of \emph{next tokens}. The words ``90\%
sure'' are themselves next tokens, chosen because they fit the text, not a measurement of how often the model
is right. Calibrated confidence has to be measured from outside, by checking many answers. Chapter 7 returns
to this.
\end{wrongidea}

\begin{checkq}
A model writes a 10-word reply, and each word is right with probability 0.95. Which statement is true?
\begin{options}
\item The reply is right with probability 0.95, because each word is.
\item The reply is right with probability $0.95^{10}$, which is less than 0.95.
\item If the reply is wrong, the cause is a bug in the code.
\item The model's reply ``I am 95\% sure'' would tell us the same number.
\end{options}
\end{checkq}

\begin{principle}{turn the task into a game the data already plays}
The strongest move in machine learning is often not a better model but a better question. ``Choose a reply
from my list'' needs labelled data and caps the system at what you imagined. ``Guess the next word'' needs no
labels and caps nothing: translation (\emph{English: waste of time. French:}), classification
(\emph{Review: waste of time. Sentiment:}) and question answering all become ``continue the text''. When you
design an AI system (an agent, a memory layer, an evaluator), ask first: what question can my data answer by
itself?
\end{principle}

So the problem is fixed for the rest of these notes:
\begin{center}
\emph{Build a function that reads $w_1,\ldots,w_{t-1}$ and returns $P(w_t = w \mid w_1,\ldots,w_{t-1})$
for every word $w$ in the vocabulary.}
\end{center}

\Summary{Writing is repeated next-word guessing, which turns every text into free, self-labelled practice,
and also lets an early mistake carry forward.}

%---------------------------------------------------------------------
\subsection*{1.3\quad Attempt: count what came next (1948 onwards)}

\Picture Keep a giant notebook. Every time you read two words, write down the word that came after them. To
guess what follows \emph{waste of}, open the notebook at the page \emph{waste of} and pick whatever was written
there most often.

\begin{breaks}
A real notebook has a page for every pair you might meet. This one only has pages for pairs it has already
read, and, as we are about to measure, that is a tiny fraction of the pairs that exist.
\end{breaks}

\Exactly This is the \emph{$n$-gram model}. Claude Shannon introduced the idea in 1948, generating
English-like text from counts of letter and word sequences, and in 1951 he measured how well \emph{people}
play the next-letter game. From the 1970s, speech-recognition groups, notably Frederick Jelinek's at IBM, built
practical systems on word trigrams:
\[
P(w_t \mid w_{t-2}, w_{t-1}) \;\approx\; \frac{\text{count}(w_{t-2}, w_{t-1}, w_t)}{\text{count}(w_{t-2}, w_{t-1})}.
\]
In words: of all the times you saw the two previous words, in what fraction did this word come next?

\Predict Take a corpus of transcribed technical talks on AI engineering (about 245,000 words), split into nine
parts. Use parts 1--8 (\cTrainWords\ words, \cVocab\ distinct words) as the notebook, and part 9
(\cTestWords\ words) as new text. Before looking at the figure: what fraction of the three-word sequences in
part 9 do you think never appear in parts 1--8? A tenth? A
quarter? Half? Write your guess down.

\begin{figure}[h]
\centering
\begin{tikzpicture}
\begin{axis}[ybar, width=9cm, height=4.6cm, bar width=18pt, ymin=0, ymax=110,
  xtick={1,2,3,4,5}, xticklabels={1 word,2 words,3 words,4 words,5 words},
  ylabel={\% never seen before}, ylabel style={font=\small}, tick label style={font=\small},
  nodes near coords, nodes near coords style={font=\footnotesize}, axis lines*=left,
  enlarge x limits=0.15]
\edef\tmp{\noexpand\addplot[fill=black!55, draw=none] coordinates {\cBars};}\tmp
\end{axis}
\end{tikzpicture}
\caption{Sparsity, measured. Sequences in part 9 that never occur in parts 1--8. Single words are
almost always known; by three words, \cUnseenThree\% are new.}
\end{figure}

Single words are almost always known (\cUnseenOne\% unseen). But \cUnseenThree\% of three-word sequences, and
\cUnseenFive\% of five-word sequences, were never seen. And our running phrase? In parts 1--8 the
page \emph{waste of} is completely blank: it occurs \cWasteOfTrain\ times. In part 9 the phrase
\tok{waste of time} occurs \cWasteOfTimeTest\ times. A counting model built from \cTrainWords\ words has nothing
to say about it.

\Sanity Why so empty? With \cVocab\ words there are about \cPossibleTriBillions\ billion possible three-word
sequences; the notebook has seen \cSeenTri\ of them, that is \cTriCoverage\% of the space. More text helps, but
the space grows as $V^n$, much faster than any library.

\Why Clever patches (smoothing, such as Kneser and Ney, 1995) move a little probability onto unseen pages, but
they cannot fix the real problem: \emph{to a counter, every word is a name and nothing more}. Having seen
\emph{waste of money} a thousand times tells it nothing about \emph{waste of time}, because it has no idea that
\emph{money} and \emph{time} are alike.

\begin{wrongidea}{1.5}
\emph{A language model looks up sentences it has already seen.} It was trained on a huge amount of text, so
storing and retrieving sounds plausible. The figure says otherwise: even inside talks on one
topic, most three-word sequences in a new part were never seen before. A model that only looked things up
would be silent most of the time. Whatever a useful model does, it has to work on text it has never seen.
\end{wrongidea}

\begin{wrongidea}{1.8}
\emph{A counting model knows that \emph{time} and \emph{money} are alike.} You know it, and the model has seen
both words many times, so it feels like it should. But the notebook stores counts of exact sequences. Nothing in
it connects the page \emph{waste of money} to the page \emph{waste of time}. That connection is what \S1.5 adds.
\end{wrongidea}

\begin{checkq}
A three-word counting model is built from parts 1--8 and is given \tok{waste of}. What does it predict next?
\begin{options}
\item \tok{time}, because \tok{waste of time} is a common English phrase.
\item Nothing: \tok{waste of} never occurs in those parts, so the count is zero.
\item \tok{time}, because it has seen \tok{waste of money} and knows \emph{money} and \emph{time} are alike.
\item \tok{time}, because \tok{waste of time} occurs \cWasteOfTimeTest\ times in part 9.
\end{options}
\end{checkq}

\begin{sidebar}
\textbf{Requirement 1.} Words must have a notion of \emph{similarity}, so that what is learned about one word
transfers to similar words.
\end{sidebar}

\Summary{Counting what came next fails because almost every sequence is new, and a counter cannot transfer
what it knows to similar words.}

%---------------------------------------------------------------------
\subsection*{1.4\quad Attempt: forget the order, count the words}

A tempting shortcut, long used in search engines and spam filters: throw the words into a bag and only count
them. A bag of \emph{waste}, \emph{of}, \emph{time} is never sparse.

\Predict Do these two customer replies give the same bag of words?
\begin{center}
\tok{not useful, very good} \qquad and \qquad \tok{very useful, not good}
\end{center}
They contain exactly the same words the same number of times, so their bags are identical (our script checks
this). They mean opposite things. Order is not decoration; it carries meaning.

\begin{wrongidea}{1.6}
\emph{If the right words are present, their order hardly matters.} Keyword search works well in daily life, so
this is easy to believe. The pair above has the right words and the opposite meaning. Order matters most for
exactly the small words, like \emph{not}, that the rules of \S1.2 tripped on.
\end{wrongidea}

\begin{checkq}
Which pair of replies would a bag-of-words model treat as identical?
\begin{options}
\item \tok{good, not cheap} and \tok{cheap, not good}
\item \tok{waste of time} and \tok{waste of money}
\item \tok{not useful} and \tok{useless}
\item None: a bag of words keeps the order of the first and last word.
\end{options}
\end{checkq}

\begin{sidebar}
\textbf{Requirement 2.} The model must know \emph{which word came where}.
\end{sidebar}

\Summary{Counting words without their order removes sparsity but also removes meaning, so order must be kept.}

%---------------------------------------------------------------------
\subsection*{1.5\quad Attempt: words as points in space (2003, 2013)}

\Picture Imagine a map of a city where every word lives at an address. \emph{Time} lives next door to
\emph{money} and \emph{effort} (things you can waste); \emph{star} lives across town. If you learn something
about one house, you have learned a little about its neighbours.

\begin{breaks}
A city map has two directions. Real word maps have hundreds, so a word can be near \emph{money} in one
direction (things you spend) and near \emph{clock} in another (things you measure). No flat map can show that.
\end{breaks}

\Exactly In the counting model, each word is effectively a \emph{one-hot} vector: a list of $V$ numbers, all zero
except a single 1 at the word's own position. Any two different one-hot vectors are exactly the same distance
apart ($\sqrt2 \approx \hOneHotDist$), so every word is an equal stranger to every other. Bengio and colleagues
(2003) replaced this with a short learned vector of $d$ numbers per word, an \emph{embedding}, and trained a
small neural network to predict the next word from the embeddings of the previous few. Because similar words
end up with similar vectors, knowledge spreads between them. Mikolov and colleagues' \emph{word2vec} (2013) made
such vectors cheap to learn from billions of words.

\Why This meets Requirement 1, and it gives us the most important idea in these notes: \emph{a word is an arrow
in space, and meaning is direction.} Chapter 2 builds it properly; Chapter 3 uses it on every page. But
Bengio's model still looked at a fixed window of the last few words. What about the word fifty words back?

\Summary{Giving each word a learned position in space lets knowledge spread between similar words, but a fixed
window still cannot reach far back.}

%---------------------------------------------------------------------
\subsection*{1.6\quad Attempt: read word by word, keep a running summary (1990 to 2014)}

\Picture Read a story one word at a time, and after each word update a single sticky note that summarises
everything so far. To guess the next word, look only at the sticky note.

\begin{breaks}
A person can reread the story; this reader cannot. Once a detail is rubbed off the note, it is gone.
\end{breaks}

\Exactly This is the \emph{recurrent neural network} (RNN; Elman, 1990; used as a language model by Mikolov and
colleagues, 2010). It keeps a hidden state $\bm{h}_t$ and updates it with each new word:
\[
\bm{h}_t = f(\bm{h}_{t-1}, \bm{x}_t), \qquad P(w_{t+1}\mid w_{\le t}) = g(\bm{h}_t).
\]
It meets Requirement 2 for free, since the note is written in reading order, and in principle it can remember
forever. In practice it breaks in three ways.

\paragraph{Break (i): one sticky note.} Whatever the sentence (three words or three thousand), it must fit into
the same fixed-size $\bm{h}_t$. Detail is squeezed out.

\paragraph{Break (ii): memories fade or explode.}
\Predict Take a single number carried on the note, and let's assume each update multiplies it by 0.9. How much
of it is left after 50 words? More or less than a quarter?

After 10 words \hFadeTen\ of it is left, after 50 words \hFadeFifty, after 100 words $\hFadeHundred$. If the
factor is 1.1, it grows to \hGrowTen, then \hGrowFifty, then $\hGrowHundred$. Real networks multiply by
matrices, not numbers, but the effect is the same, and it hits learning as hard as remembering: the signal that
should teach the network ``the word fifty steps back mattered'' has to travel back through the same fifty
multiplications (Bengio, Simard and Frasconi, 1994). The LSTM (Hochreiter and Schmidhuber, 1997) added
\emph{gates}, learned switches that let the note keep some numbers untouched, which eased this a lot but did not
remove breaks (i) or (iii).

\paragraph{Break (iii): it cannot be done in parallel.} You cannot write note 1000 before note 999, and you
cannot write 999 before 998. A sentence of \hT\ tokens is \hT\ steps that must happen one after another. A GPU
is a machine with thousands of small workers; give it a chain like this and almost all of them wait.

\begin{sidebar}
\textbf{Requirement 3.} Any word must be able to reach any earlier word \emph{directly}, not through a chain of
summaries.\\
\textbf{Requirement 4.} The work for all positions must be doable \emph{at the same time}.
\end{sidebar}

\begin{checkq}
An RNN carries a signal that is multiplied by 0.9 at every word. Which statement is true?
\begin{options}
\item After 50 words, the signal is about half its size.
\item After 50 words, less than 1\% of the signal is left, and the lesson ``this early word mattered'' fades the same way.
\item An LSTM's gates make the RNN process all words at the same time.
\item The fading only affects remembering, not learning.
\end{options}
\end{checkq}

\Summary{An RNN reads in order, but squeezes everything into one note, lets old signals fade, and cannot work on
all positions at once.}

%---------------------------------------------------------------------
\subsection*{1.7\quad The turn: let the reader look back (2014 to 2017)}

In 2014, neural translation systems (Sutskever, Vinyals and Le; Cho and colleagues) read an English sentence
with one RNN into a single vector and wrote the French with another. Break (i) was plain to see: long sentences
came out worse. Bahdanau, Cho and Bengio (2014) proposed a fix: when writing each French word, let the writer
\emph{look back at every English word} and decide how much each one matters right now. They called these weights
\emph{attention}. Requirement 3 was met. But there was still an RNN underneath, so Requirement 4 was not.

In 2017, Vaswani and colleagues asked: \emph{what if attention is all you need?} Drop the RNN entirely. Let every
word look at every earlier word directly, and do it for all words at once. Figure~\ref{fig:rnn-vs-attn} shows
the difference. Their base translation model trained in about twelve hours on eight GPUs. That architecture is
the \emph{transformer}; GPT (2018), GPT-2 (2019) and every model in Chapter 10 are transformers trained on the
next-word game of \S1.2.

\begin{figure}[h]
\centering
\begin{tikzpicture}[font=\small, w/.style={draw, circle, minimum size=0.8cm, inner sep=0pt, line width=0.7pt}, x=1.2cm]
\begin{scope}
\foreach \i/\t in {1/w_1,2/w_2,3/w_3,4/w_4,5/w_5} \node[w] (r\i) at (1.15*\i,0) {$\t$};
\foreach \i [evaluate=\i as \j using int(\i+1)] in {1,...,4} \draw[-{Stealth}] (r\i) -- (r\j);
\node[align=center] at (3.45,-1.05) {RNN: $w_5$ reaches $w_1$ through 4 hops;\\5 steps, one after another};
\end{scope}
\begin{scope}[xshift=7.6cm]
\foreach \i/\t in {1/w_1,2/w_2,3/w_3,4/w_4,5/w_5} \node[w] (a\i) at (1.15*\i,0) {$\t$};
\foreach \i in {1,...,4} \foreach \j in {2,...,5} {
  \ifnum\i<\j \draw[-{Stealth}, black!70, line width=0.6pt] (a\j) to[bend right=40] (a\i); \fi }
\node[align=center] at (3.45,-1.05) {Attention: every word reaches every earlier\\word in 1 hop; all pairs computed at once};
\end{scope}
\end{tikzpicture}
\caption{The two shapes of computation. For a sentence of \hT\ tokens, the RNN needs \hT\ dependent steps and
\hHops\ hops from last to first; attention computes all \hTT\ pairings in one parallel step and connects any two
words in one hop.}
\label{fig:rnn-vs-attn}
\end{figure}

\begin{wrongidea}{1.7}
\emph{Transformers replaced RNNs mainly because they are more accurate.} Better results are the headline of
every paper, so this is the natural reading. Look at the two breaks the transformer removed: the long chain
between distant words (Requirement 3) and the one-step-at-a-time work (Requirement 4). The second one decided
the history: a model that uses every GPU worker at once can be trained on far more text in the same time, and
the accuracy followed from the scale.
\end{wrongidea}

\begin{center}
\small
\begin{tabular}{@{}l l l@{}}
\toprule
Attempt & What it fixed & Where it broke \\
\midrule
Rules & nothing to learn & endless paraphrases, negation\\
$n$-gram counting (1948 on) & learns from text, no labels & \cUnseenThree\% of 3-word sequences unseen\\
Bag of words & no sparsity & order lost: opposite meanings, same bag\\
Embeddings (2003, 2013) & similarity between words & fixed window of recent words\\
RNN, LSTM (1990, 1997) & order; unlimited history in principle & one note; fading; strictly sequential\\
RNN + attention (2014) & direct look-back & still sequential underneath\\
Transformer (2017) & direct look-back, all at once & \emph{order is lost again}: see \S3.11\\
\bottomrule
\end{tabular}
\end{center}

\Summary{Attention gave every word a direct line to every earlier word, and dropping the RNN let all of that run
at once, which is what made training at scale possible.}

%---------------------------------------------------------------------
\subsection*{1.8\quad A real case: Google Translate, 2016}

For a decade, Google Translate worked phrase by phrase: cut the sentence into short chunks, look each chunk up in
a huge table of translations counted from text, then glue and reorder the pieces. That is the counting notebook
of \S1.3, applied to translation, and it had the same weakness: anything outside a chunk was invisible, so whole
sentences came out stiff. In 2016 Google replaced it with a neural system (Wu and colleagues, 2016): LSTMs that
read the whole sentence, with attention letting the writer look back at every source word. That is \S1.6 plus
\S1.7, about a year before the transformer. A year later the transformer paper reported better translation
quality than such systems at a fraction of the training cost.

Tell me: which of the four requirements does each step in this story meet? The phrase table meets none of 1, 3
and 4. The 2016 system meets 1, 2 and 3. The transformer meets 1, 3 and 4, and loses 2 for a while (Chapter 4).

\Summary{A product used by millions went through the same sequence of fixes as these notes, for the same
reasons.}

%---------------------------------------------------------------------
\subsection*{1.9\quad The whole machine at a glance}

Before we open the box, here is the whole forward pass, the route every batch of text takes from characters to
probabilities. Do not try to understand each stage yet; look at the shape of the road.

\Picture An assembly line. Sentences come in at one end as text. They are cut into pieces (tokens), each piece
gets its starting arrow (embedding), and then the arrows pass through a row of identical workstations. At each
station two things happen: the tokens \emph{talk to each other} (attention), then each token \emph{thinks on its
own} about what it heard (the MLP). At the end of the line, every token's arrow is turned into a guess about the
word that comes after it.

\begin{breaks}
On an assembly line each item passes each station once and the stations never look back. Here every station
sees all the tokens of a sentence together, and what one token becomes depends on the others.
\end{breaks}

\paragraph{Drawing the transformer, block by block.} Here is the whole machine as a frame of empty blocks.
Every chapter from now on fills in one or two of them, and opens with this same picture, so you always know
where you are: solid blocks are built, the highlighted block is the one that chapter builds, and dashed blocks
are still ahead. For now only the top is filled in: text comes in.

\machine{1}{The frame we will fill. Text enters at the top; each dashed block is one chapter's job. The middle
part, inside the large rounded box, is one \emph{layer}; a real model stacks $L$ of them. At the bottom, the
model gives a score to every word it knows, of shape $B\times T\times V$: a guess at \emph{every} position of
\emph{every} sentence in one pass, not only at the last one.}

\begin{center}% AUTO-GENERATED by code/figlib.py helpers -- do not edit
\begin{tikzpicture}[x=1.7cm]
\draw[black!45, fill=white] (0.000,0.000) rectangle (0.780,-0.640);
\node[font=\small, text=black] at (0.390,-0.320) {\texttt{waste}};
\draw[black!45, fill=white] (0.780,0.000) rectangle (1.560,-0.640);
\node[font=\small, text=black] at (1.170,-0.320) {\texttt{of}};
\draw[black!45, fill=white] (1.560,0.000) rectangle (2.340,-0.640);
\node[font=\small, text=black] at (1.950,-0.320) {\texttt{time}};
\draw[black!45, fill=white] (0.000,-0.640) rectangle (0.780,-1.279);
\node[font=\small, text=black] at (0.390,-0.959) {\texttt{not}};
\draw[black!45, fill=white] (0.780,-0.640) rectangle (1.560,-1.279);
\node[font=\small, text=black] at (1.170,-0.959) {\texttt{very}};
\draw[black!45, fill=white] (1.560,-0.640) rectangle (2.340,-1.279);
\node[font=\small, text=black] at (1.950,-0.959) {\texttt{useful}};
\draw[line width=0.9pt] (0.000,0.000) rectangle (2.340,-1.279);
\node[anchor=east, font=\footnotesize\ttfamily] at (-0.080,-0.320) {sentence 1};
\node[anchor=east, font=\footnotesize\ttfamily] at (-0.080,-0.959) {sentence 2};
\node[anchor=south, font=\footnotesize\ttfamily] at (0.390,0.020) {1};
\node[anchor=south, font=\footnotesize\ttfamily] at (1.170,0.020) {2};
\node[anchor=south, font=\footnotesize\ttfamily] at (1.950,0.020) {3};
\node[anchor=south, font=\normalsize] at (1.170,0.420) {what the model reads};
\node[anchor=north, font=\footnotesize, text=black!75] at (1.170,-1.359) {$B\times T=2\times3$};
\end{tikzpicture}
\qquad
\begin{tikzpicture}[x=1.7cm]
\draw[black!45, fill=white] (0.000,0.000) rectangle (0.780,-0.640);
\node[font=\small, text=black] at (0.390,-0.320) {\texttt{of}};
\draw[black!45, fill=white] (0.780,0.000) rectangle (1.560,-0.640);
\node[font=\small, text=black] at (1.170,-0.320) {\texttt{time}};
\draw[black!45, fill=white] (1.560,0.000) rectangle (2.340,-0.640);
\node[font=\small, text=black] at (1.950,-0.320) {\texttt{Sorry}};
\draw[black!45, fill=white] (0.000,-0.640) rectangle (0.780,-1.279);
\node[font=\small, text=black] at (0.390,-0.959) {\texttt{very}};
\draw[black!45, fill=white] (0.780,-0.640) rectangle (1.560,-1.279);
\node[font=\small, text=black] at (1.170,-0.959) {\texttt{useful}};
\draw[black!45, fill=white] (1.560,-0.640) rectangle (2.340,-1.279);
\node[font=\small, text=black] at (1.950,-0.959) {\texttt{.}};
\draw[line width=0.9pt] (0.000,0.000) rectangle (2.340,-1.279);
\node[anchor=east, font=\footnotesize\ttfamily] at (-0.080,-0.320) {sentence 1};
\node[anchor=east, font=\footnotesize\ttfamily] at (-0.080,-0.959) {sentence 2};
\node[anchor=south, font=\footnotesize\ttfamily] at (0.390,0.020) {1};
\node[anchor=south, font=\footnotesize\ttfamily] at (1.170,0.020) {2};
\node[anchor=south, font=\footnotesize\ttfamily] at (1.950,0.020) {3};
\node[anchor=south, font=\normalsize] at (1.170,0.420) {the word it must guess next};
\node[anchor=north, font=\footnotesize, text=black!75] at (1.170,-1.359) {$B\times T=2\times3$ guesses};
\end{tikzpicture}
\end{center}

Look at the two tables above. That is the quiet reason transformers train so fast. A sentence of $T$ tokens is not one exercise
but $T$ exercises: after \tok{waste}, guess \tok{of}; after \tok{waste of}, guess \tok{time}; after \tok{waste
of time}, guess \tok{Sorry}. The transformer answers all of them in the same pass, for all $B$ sentences in the
batch. How can it do that without cheating, without the guess for position 2 peeking at the answer in position
3? That is one of the things Chapter 3 has to explain.

\begin{checkq}
A batch holds $B=8$ sentences of $T=16$ tokens each. How many next-word guesses does one forward pass produce?
\begin{options}
\item 8, one per sentence.
\item 16, one per position.
\item 128, one per position of every sentence.
\item It depends on how long the reply is.
\end{options}
\end{checkq}

\Summary{One pass turns $B$ sentences of $T$ tokens into $B\times T$ next-word guesses, through tokens,
embeddings, $L$ talk-then-think stations and a final scoring step.}

%---------------------------------------------------------------------
\begin{lookingback}
These look back at Chapter 0, the notation page.
(1) What do $B$, $T$ and $d$ stand for, in words? (2) A matrix has shape $3\times4$. How many rows does it have,
and what does each row hold in these notes? (3) Which symbol will stand for the number of learned parameters?
\end{lookingback}

\subsection*{Exercises}

\begin{exercise}
A notebook model looks at the previous 4 words, and the vocabulary has 50{,}000 words. How many pages could the
notebook have? Is that more or less than the number of words ever written?
\end{exercise}

\begin{exercise}
Using the bag-of-words idea, find two more replies a support team might receive that have identical bags but
opposite meanings.
\end{exercise}

\begin{exercise}
An RNN multiplies a signal by 0.95 per word. After how many words is less than 1\% of it left?
\end{exercise}

\begin{exercise}
Write the chain rule for the reply \tok{Sorry , this has not been useful} given \tok{waste of time}: how many
factors are there, and what does the third factor condition on?
\end{exercise}

\subsection*{Answers to Chapter 1}
\small
\textbf{Check yourself 1.1.} (b). A chain of guesses multiplies, so the whole is less likely than each part:
$0.95^{10}\approx\nOneTenRight$. (a) treats the whole reply as one guess. (c) is the wrong idea M1.3: wrong words come
from the guessing process itself. (d) is M1.2: ``I am 95\% sure'' is more text, not a measurement.

\textbf{Check yourself 1.2.} (b). The page \emph{waste of} is blank: \cWasteOfTrain\ occurrences. (a) is M1.5:
the model does not look things up in ``English''; it only has its notebook. (c) is M1.8: counts do not know
which words are alike. (d) uses data the model never trained on, the same leakage your fraud model must avoid.

\textbf{Check yourself 1.3.} (a). Same words, same counts, opposite meanings. (b) differ in one word. (c) have
different words. (d) is M1.6 in disguise: a bag keeps no order at all.

\textbf{Check yourself 1.4.} (b). $0.9^{50}=\hFadeFifty$. (a) underestimates how fast repeated multiplication
shrinks a number. (c) is wrong: gates help memory, not parallelism (M1.7 is the belief that accuracy, not
parallelism, was the deciding issue). (d) misses that learning signals travel back through the same steps.

\textbf{Check yourself 1.5.} (c). $8\times16=128$: every position of every sentence predicts its next token. (a)
and (b) count one dimension only. (d) confuses training on given text with writing a new reply.

\textbf{Looking back.} (1) Batch size, tokens per sentence, numbers per token. (2) Three rows, one per token.
(3) $N$.

\textbf{Exercise 1.1.} $50{,}000^4 = \hExPages$ possible four-word contexts; far more than all words ever
written, so almost every page stays empty.

\textbf{Exercise 1.2.} For example \tok{good, not cheap} (good but expensive) versus \tok{cheap, not good}.

\textbf{Exercise 1.3.} We need $0.95^n < 0.01$, that is $n > \ln 0.01/\ln 0.95 \approx \hExN$, so after
\hExNceil\ words.

\textbf{Exercise 1.4.} Seven factors, one per reply token. The third is $P(\tok{this}\mid\tok{waste of time Sorry ,})$:
it conditions on the message and the two reply tokens before it.
\normalsize

\subsection*{What you now know}
\begin{itemize}
\item Why writing a reply one word at a time beats rules and fixed reply lists, and why that game labels itself.
\item How the chain rule turns a sentence into a product of next-word probabilities, and why errors compound.
\item Why counting, bag of words, word vectors, RNNs and RNN-plus-attention each fell short, measured where
possible, and the four requirements those failures left behind.
\item The shape of the whole forward pass, and why one pass gives $B\times T$ guesses.
\end{itemize}
